\section{Data}
\label{sec:data}
\subsection{National Registry of Disappeared and Not Located Persons}

The National Registry of Missing and Unlocated Persons (\textit{Registro Nacional de Personas Desaparecidas y No Localizadas}, RNPDNO) is Mexico's centralized registry of disappearance cases, established under the 2017 General Law on Disappearances \citep{ley_general_2017} and administered by the National Search Commission (\textit{Comisión Nacional de Búsqueda}, CNB). The CNB publishes aggregate statistics and visualizations through a public dashboard \citep{cnb_rnpdno} but does not release machine-readable, municipality-level data files. The dataset used in this study was constructed by web-scraping, geocoding, and processing the dashboard records into four structured CSV files, one per outcome status \citep{escobar_rnpdno_2026}.

The registry records new disappearance cases each month, not a running total of all cases since the registry was created. During the \nmonths{}-month study period \studyperiod{}, the registry recorded 263,402 person-registrations, of which 89,631 were classified as Not Located, 157,990 as Located Alive, and 15,438 as Located Dead (Table~\ref{tab:dataset}). Municipality-month combinations with no recorded cases are treated as zero registrations, yielding a balanced panel of \nmunis{} municipalities observed over \nmonths{} months and four outcome categories. Records with unresolvable municipality identifiers (12,206 persons, 4.6\% of Total registrations) are excluded from spatial analyses, potentially biasing hotspot maps toward municipalities with more complete geocoding. Each file additionally provides sex-disaggregated counts (male, female, undefined).
\begin{tabular}{lrrrrrrrr}
  \toprule
  & & & & \multicolumn{4}{c}{\textbf{Valid persons (by sex)}} \\
  \cmidrule(lr){5-8}
  \textbf{Status} & \textbf{Valid rows} & \textbf{Sentinel rows} & \textbf{Valid persons} & \textbf{Male} & \textbf{Female} & \textbf{Undef.} & \textbf{\%Male} \\
  \midrule
  Total          & 55,524 & 1,803 & 262,814 & 160,857 & 101,630 & 327 & 61.2\% \\
  Not Located    & 33,384 & 1,256 &  89,850 &  69,871 &  19,843 & 136 & 77.8\% \\
  Located Alive  & 36,521 &   837 & 157,688 &  77,856 &  79,686 & 146 & 49.4\% \\
  Located Dead   &  9,643 &   240 &  15,343 &  13,175 &   2,123 &  45 & 85.9\% \\
  \bottomrule
  \multicolumn{8}{l}{\footnotesize Temporal coverage: January 2015 -- December 2025. Spatial analysis covers 2,478 municipalities (INEGI 2024).} \\
  \multicolumn{8}{l}{\footnotesize Sentinel geocodes (cvegeo ending 998/999 or = 99998) excluded. \%Male computed over valid persons.} \\
  \multicolumn{8}{l}{\footnotesize Located Alive is the only status with near-equal sex distribution (49.4\% male). Located Dead is the most male-skewed (85.9\%).} \\
\end{tabular}


\subsection{Geographic Framework}

We use the INEGI Marco Geoestad\'{i}stico 2025 \citep{inegi_ageeml}, which defines \nmunis{} municipalities in projected coordinates (EPSG:6372, Mexico ITRF2008/LCC).
The five-digit cvegeo code serves as the spatial join key across all datasets; all \nmunis{} municipalities match one-to-one between the geoparquet and the panel dataset (100\% match rate, zero merge failures).
\subsection{Regional Classification}
Municipalities are classified into five macro-regions following the CESOP scheme \citep{mexicoencifras}: Norte, Norte-Occidente, Centro-Norte, Centro, and Sur.
The sixth region, the Baj\'{i}o industrial corridor, is defined independently of the CESOP classification.
\citet{unger_2014} characterized the Baj\'{i}o as a set of economic corridors organized along six highway axes, encompassing over 200 municipalities.
We adopt a narrower operationalization focused on the urbanized industrial core: the 29 municipalities comprising the 10 metropolitan areas and 4 conurbations identified by \citet{medina_fernandez_2023} from the \emph{Sistema Urbano Nacional} 2018 catalog \citep{conapo_sun_2018}, spanning 14 urban entities across 5 states (Aguascalientes, Guanajuato, Michoac\'{a}n, Quer\'{e}taro, and San Luis Potos\'{i}).
This narrower definition isolates the municipalities where industrial activity, population density, and cartel territorial competition converge. The corridor is not mutually exclusive with the five-region scheme; its municipalities also belong to their respective macro-regions.
Section~\ref{sec:trend} reports trend slopes for all six units.